\documentclass[12pt]{article}
\usepackage[margin=1in]{geometry}
\usepackage{hyperref}
\usepackage{booktabs}
\usepackage{parskip}
\usepackage{enumitem}
\usepackage{xcolor}
\usepackage{mdframed}
\usepackage{titlesec}

\hypersetup{colorlinks=true, linkcolor=blue, urlcolor=blue}

\titleformat{\section}{\large\bfseries}{}{0em}{}[\titlerule]
\titleformat{\subsection}{\normalsize\bfseries}{}{0em}{}

\title{\textbf{Mapping \& Analyzing Participation}\\
\large Suggestions and Guidance for Future Research Teams}
\author{E483/E583 Information Visualization --- Spring 2026\\
Indiana University Luddy School of Informatics, Computing, and Engineering}
\date{April 2026}

\begin{document}
\maketitle
\tableofcontents
\newpage

% -------------------------------------------------------------------------
\section{Project Overview}

This document is written for researchers and development teams who continue
work on the \textbf{Mapping and Analyzing Participation} project after the
Spring 2026 semester. The project was commissioned by Daniel F.~Bassill of the
\textbf{Tutor/Mentor Connection} (T/MC, 1993--2011) and its successor
\textbf{Tutor/Mentor Institute LLC} (TMI, 2011--present) to analyze and
visualize attendance records from 42 Tutor/Mentor Leadership and Networking
Conferences held in Chicago between May 1994 and November 2015.
When referencing the client organization, use T/MC for events prior to 2011
and Tutor/Mentor Institute LLC for anything after; Bassill uses both names
throughout his published materials.

The primary dataset is:
\begin{center}
\texttt{ProjectFiles/All\_Conferences\_Data\_Cleaned-2026 - Sheet1.csv}
\end{center}
It contains 6,410 participation records across 18 columns, covering roughly
2,161 unique organizations and 21 years of network-building activity in
Chicago's high-poverty neighborhoods.

The codebase is a \textbf{Django} web application with a \texttt{Pandas}/
\texttt{Plotly}/\texttt{NetworkX} analytics layer. All source code lives
in the \texttt{conferences/} app. Docker support is included via
\texttt{Dockerfile} and \texttt{docker-compose.yml}.

% -------------------------------------------------------------------------
\section{Understanding the Data Before You Touch It}

\subsection{The SNA Category System}

Every attendance record carries an \texttt{sna\_category} field drawn from a
controlled vocabulary of eleven canonical types:

\begin{center}
\begin{tabular}{ll}
\toprule
Category & What it represents \\
\midrule
Program      & Tutoring/mentoring program providers \\
Business     & Corporate sponsors and private-sector partners \\
College      & Universities and community colleges \\
K-12 School  & Public and private K-12 institutions \\
Government   & Municipal, county, state, and federal agencies \\
Faith        & Faith-based organizations \\
Foundation   & Grant-making foundations and philanthropies \\
Intermediary & Capacity-building and coalition organizations \\
T/MC         & Tutor/Mentor Connection itself \\
Resource     & Consultants and support services \\
Media        & Journalists, reporters, and media organizations \\
Other        & Records that could not be classified \\
\bottomrule
\end{tabular}
\end{center}

The original column (\texttt{sna\_category\_original}) contained hundreds of
inconsistent spellings. The normalization mapping lives in
\texttt{conferences/sna\_normalize.py}. \textbf{Consult Daniel Bassill
before adding or reclassifying categories}---the vocabulary is intentional
and affects downstream network analysis.

Bassill has explicitly requested a \textbf{12th category: Media}
(journalists and media organizations). This category is intentionally thin
in the current dataset, and that sparsity is itself meaningful---building
media coverage is a core conference goal, and its absence in the data
reflects a real gap in outreach.

\subsection{Organization Name Cleaning}

The \texttt{organization} field (cleaned) is the authoritative identifier for
network analysis. The \texttt{organization\_original} field preserves the raw
value as entered at each conference. Never use the original field as a join
key---variant spellings of the same organization appear across years.

The cleaning was performed manually by the client in the 2026 CSV revision.
If the dataset is extended with new records, the same cleaning step must be
applied before loading data into Django.

\subsection{Known Data Quality Issues}

\begin{itemize}[leftmargin=1.5em]
  \item A small number of records have a blank \texttt{organization} field.
        These are excluded from all network exports and most analytics.
  \item Records classified as \texttt{Other} may be reclassifiable with
        additional domain knowledge from the client.
  \item The \texttt{active\_inactive} column is populated inconsistently and
        was not used in the Spring 2026 analysis. It may carry useful
        longitudinal signal for organizations that later closed.
  \item Address and ZIP code fields are sparsely populated for earlier years
        (pre-2000). Geographic analysis is more reliable for 2000 onward.
  \item \textbf{Organization names change over time.} An organization that
        attended in the 1990s may appear under a different name in later
        years---sometimes a third name by the 2010s. The cleaned
        \texttt{organization} field represents the name as used at that
        conference; cross-decade queries may miss the same entity under a
        variant name.
  \item \textbf{Individuals moved between organizations.} A person who
        attended four conferences through three different employers will appear
        as three separate organizations each with one attendee, not as one
        person with four appearances. The first/last name fields can be used
        to surface this pattern without exposing names publicly.
\end{itemize}

% -------------------------------------------------------------------------
\section{Technical Infrastructure}

\subsection{Running the Application}

\textbf{With Docker (recommended):}
\begin{verbatim}
docker-compose up --build
\end{verbatim}
The app will be available at \texttt{http://localhost:8000}.

\textbf{Without Docker:}
\begin{verbatim}
pip install -r requirements.txt
python manage.py migrate
python manage.py load_data   # loads the CSV into SQLite
python manage.py runserver
\end{verbatim}

\subsection{Loading Data}

The management command \texttt{conferences/management/commands/load\_data.py}
reads the CSV, applies SNA normalization (via \texttt{sna\_normalize.py}),
and populates the \texttt{Attendance} table. Re-running it with a revised CSV
will not duplicate records because it clears and reloads the table. If you
extend the dataset, verify the column headers match those expected by the
loader before running.

\subsection{Database}

The application ships with SQLite (\texttt{db.sqlite3}). For a production
deployment with concurrent users, migrate to PostgreSQL. The Django ORM
abstracts this entirely---update \texttt{DATABASES} in
\texttt{MappingAndAnalyzingParticipation/settings.py} and add
\texttt{psycopg2-binary} to \texttt{requirements.txt}.

\subsection{Environment Variables}

Sensitive settings are stored in \texttt{.env} (not committed to git).
A \texttt{SECRET\_KEY} must be set before deployment. See \texttt{.env}
for the variable names expected.

% -------------------------------------------------------------------------
\section{Network Analysis Guidance}

\subsection{Available Exports}

The application exposes four CSV export endpoints under \texttt{/export/}:

\begin{center}
\begin{tabular}{lp{8cm}}
\toprule
Endpoint & Contents \\
\midrule
\texttt{/export/nodes/} & One row per organization with sector, event count, and record count \\
\texttt{/export/edges/org-event/} & Bipartite edges: organization $\leftrightarrow$ conference event \\
\texttt{/export/edges/org-org/} & Co-attendance edges with shared-conference weights \\
\texttt{/export/edges/event-sector-org/} & Three-layer (event $\to$ sector $\to$ org) edge list \\
\bottomrule
\end{tabular}
\end{center}

All endpoints accept \texttt{?year=}, \texttt{?sector=}, and
\texttt{?state=} query parameters for filtered exports.

\subsection{Importing into Kumu.io}

The application exposes JSON endpoints (e.g.\ \texttt{/export/edges/org-event/})
that Kumu can link to directly. This is the preferred approach because the
map updates automatically when the underlying data changes:

\begin{enumerate}[leftmargin=1.5em]
  \item Copy the JSON export URL from the Network Maps page on the running app.
  \item In Kumu, choose \textbf{Import} $\to$ \textbf{Link map to remote JSON}
        and paste the URL.
  \item Use \textbf{Import} $\to$ \textbf{Re-import} to force a refresh without
        recreating the map from scratch.
\end{enumerate}

Alternatively, for a fully static import via CSV:

\begin{enumerate}[leftmargin=1.5em]
  \item Download \texttt{nodes.csv} and \texttt{edges\_org\_org.csv}.
  \item In Kumu, create a new map and use \textbf{Import} $\to$ \textbf{CSV}.
  \item Map \texttt{Id}/\texttt{Label} to Kumu's element fields and
        \texttt{SNA\_Category} as the element type for color coding.
  \item Use \texttt{Event\_Count} as the size field to weight nodes by
        participation history.
\end{enumerate}

When embedding network maps in the Django app, \textbf{link to the live Kumu
map URL} rather than rendering a static PNG. A linked map lets visitors open
a full interactive view, which is significantly easier to explore than a small
embedded image.

\subsection{Importing into Gephi}

\begin{enumerate}[leftmargin=1.5em]
  \item Download \texttt{nodes.csv} and \texttt{edges\_org\_org.csv}.
  \item In Gephi, use \textbf{File} $\to$ \textbf{Import Spreadsheet}.
        Import nodes first, then edges.
  \item Set \texttt{Source}/\texttt{Target} as edge columns and
        \texttt{Weight} as the edge weight.
  \item Run \textbf{ForceAtlas 2} layout. Set node size to
        \texttt{Event\_Count} and color by \texttt{SNA\_Category}.
  \item Recommended analysis modules: modularity (community detection),
        betweenness centrality, and eigenvector centrality to identify
        bridge organizations.
\end{enumerate}

\subsection{Suggested Network Research Questions}

\begin{itemize}[leftmargin=1.5em]
  \item Which organizations served as consistent \emph{bridges} between
        sectors (high betweenness centrality) across the full 21-year span?
  \item Did the network become more or less cross-sector connected over time?
        Run year-sliced graphs and compare modularity scores.
  \item Are there organizations that attended many early conferences but
        disappeared from later ones? Do they correlate with the
        \texttt{active\_inactive} field?
  \item Which sector pairs co-attended most frequently? This can inform
        partnership recommendations for T/MC's current outreach.
\end{itemize}

% -------------------------------------------------------------------------
\section{Recommended Next Steps}

\subsection{High Priority}

\begin{enumerate}[leftmargin=1.5em]
  \item \textbf{Geographic mapping.}
        ZIP codes are present for a substantial portion of records. Joining
        to a Chicago neighborhood shapefile (available from the Chicago
        Data Portal) would allow choropleth maps showing which neighborhoods
        were most represented and which were underserved---directly relevant
        to Bassill's outreach mission.

  \item \textbf{Organization URL enrichment completion.}
        The script \texttt{find\_org\_websites.py} located URLs for most
        organizations and saved results to \texttt{org\_websites.csv}. These
        URLs have not yet been joined to the Django database. Adding a
        \texttt{website} field to the \texttt{Attendance} model and linking
        via \texttt{organization} name would enable clickable org profiles
        in the dashboard.

  \item \textbf{Manual review of \texttt{unverified} and
        \texttt{not found} URL results.}
        Roughly a quarter of organizations in \texttt{org\_websites.csv}
        have status \texttt{unverified} or \texttt{not found}. Cross-referencing
        with the Tutor/Mentor Exchange directory at
        \href{https://tutormentorexchange.net}{tutormentorexchange.net}
        would improve coverage. Note: always use the \texttt{.net} domain---
        \texttt{.com} does not resolve.

  \item \textbf{OHATS data integration.}
        The OHATS (Organization Health Assessment Tool \& Survey) dataset is
        partially cleaned and ready for a future team to pick up. It is a
        strong candidate for the next semester project because it is
        data-visualization-heavy and connects directly to Bassill's four
        strategic pillars (see the ``Strategy'' concept map on
        \href{https://tutormentorexchange.net}{tutormentorexchange.net} under
        Tutor/Mentor Connection $\to$ Strategy). Future teams should read the
        OHATS section of the project handbook before beginning.

  \item \textbf{Add a ``Media'' SNA category.}
        Update \texttt{conferences/sna\_normalize.py} to map media-related
        original values to a new \texttt{Media} category and remove
        \texttt{Resource} entries that properly belong there. Confirm
        the mapping with Bassill before committing.
\end{enumerate}

\subsection{Medium Priority}

\begin{enumerate}[leftmargin=1.5em]
  \item \textbf{Longitudinal trend analysis.}
        The participation frequency chart currently shows absolute attendance
        counts. Normalizing by the number of organizations \emph{known to
        exist} in each year (using external directories) would convert raw
        counts to participation rates---a more meaningful metric for
        researchers.

  \item \textbf{Individual-level analysis.}
        The dataset includes first and last names, enabling tracking of
        individual attendees across conferences. Identifying individuals who
        attended the most events, changed organizations over time, or bridged
        sectors as individuals (not just organizationally) would add a
        person-centered perspective to the network view. Be mindful of
        privacy: individual-level data should not be published without
        Bassill's explicit approval.

  \item \textbf{Interactive network visualization in-browser.}
        The three current network maps are static PNGs rendered server-side
        by Matplotlib. Replacing them with a D3.js or Sigma.js force-directed
        graph would allow users to hover, filter, and explore nodes
        interactively without downloading Gephi or Kumu.

  \item \textbf{Dashboard usage guide.}
        The category/year/conference filter controls are not self-explanatory
        to a first-time visitor. Adding a brief on-page explanation or tooltip
        layer describing what each filter does---and how the top-10 view can
        surface duplicate or miscategorized organizations---would significantly
        improve usability.

  \item \textbf{Link to fall 2025 second group's work.}
        A second team also analyzed the conference data in fall 2025 (in
        addition to Team K). Both teams' work should be linked under the
        ``More'' dropdown in the Network Maps section of the app so the full
        lineage of this project is visible to future researchers.

  \item \textbf{Cross-reference the printed conference directories.}
        Bassill produced printed directories of invited organizations from
        1994 to 2003 (available as PDFs). Comparing those invitation lists
        to the actual attendance records would reveal who was invited but
        never attended---a meaningful gap for outreach analysis.

  \item \textbf{Public deployment.}
        The Docker configuration is ready. Deploying to a cloud provider
        (Render, Railway, or a university-hosted server) with a persistent
        PostgreSQL database would make the tool accessible to Bassill and the
        broader research community without requiring a local Python environment.
\end{enumerate}

\subsection{Low Priority / Long-Term}

\begin{enumerate}[leftmargin=1.5em]
  \item \textbf{Comparative analysis.}
        The network structure of this Chicago conference series could be
        compared to similar civic networking efforts in other cities using
        published SNA datasets, positioning this work within the broader
        nonprofit network literature.

  \item \textbf{IU School of Philanthropy collaboration.}
        Brad Fulton (IU Lilly Family School of Philanthropy) leads Project 990,
        which maps nonprofit organizations using IRS data. Connecting that
        geographic nonprofit database to the conference participation data
        could reveal which types of organizations are most underrepresented
        at conferences relative to their presence in a given community.
        Fulton has connected with Bassill on LinkedIn and has expressed
        interest; reaching out through the IU Data Science Club or an
        independent study arrangement is a natural entry point.

  \item \textbf{Donor pipeline layer.}
        Bassill described a concept for aggregating donor-flow data on top of
        the geographic nonprofit map: tracking how many visitors arrived at
        a nonprofit's donation portal through a shared platform, broken down
        by ZIP code and time period. This is a longer-term product idea that
        would require partnerships with nonprofit payment processors and is
        well suited to a capstone or MBA-level project team.

  \item \textbf{Visualization competition outreach.}
        The IU Luddy hackathon and Tableau's annual Iron Viz competition are
        both viable venues for surfacing this dataset to a wider student
        audience. Bassill has also pointed to existing websites that aggregate
        student visualization work as models for a permanent archive of
        conference-data visualizations done by successive teams.
\end{enumerate}

% -------------------------------------------------------------------------
\section{Working With the Client}

Daniel F.~Bassill (\texttt{tutormentor1@gmail.com}) is the domain expert
for all questions about the data's meaning, the SNA category definitions,
and organizational context. He has deep knowledge of Chicago's tutoring
and mentoring ecosystem that no documentation can fully substitute.

When writing about the client organization, note that \textbf{Tutor/Mentor
Connection} (T/MC) is the historically correct name for the 1993--2011 era;
\textbf{Tutor/Mentor Institute LLC} (TMI) is the current legal entity
(2011--present). Bassill uses both names in his published materials and PDF
essays. Do not refer to T/MC as the current organization---doing so
understates the institutional transition. Blog resources are at
\href{https://tutormentor.blogspot.com}{tutormentor.blogspot.com}; tags for
``network analysis,'' ``Kumu,'' and ``tipping point'' are particularly
useful starting points for new team members.

Key things to confirm with him before major analytical decisions:

\begin{itemize}[leftmargin=1.5em]
  \item Any reclassification of \texttt{Other} records into named categories.
  \item Any changes to organization name normalization.
  \item Whether individual-level data (names) can appear in public-facing views.
  \item The intended audience for any new visualizations---his four
        stakeholder groups (T/MC, nonprofits, researchers, general public)
        have different technical literacy levels.
\end{itemize}

% -------------------------------------------------------------------------
\section{Repository and File Reference}

\begin{center}
\begin{tabular}{lp{8.5cm}}
\toprule
Path & Purpose \\
\midrule
\texttt{conferences/models.py} & Django ORM model for attendance records \\
\texttt{conferences/views.py} & All dashboard, chart, export, and map views \\
\texttt{conferences/sna\_normalize.py} & SNA category normalization mapping \\
\texttt{conferences/management/commands/load\_data.py} & CSV ingestion command \\
\texttt{ProjectFiles/*.csv} & Source data and org URL output \\
\texttt{find\_org\_websites.py} & Web crawler for organization URLs \\
\texttt{generate\_network\_maps.py} & Standalone script for batch PNG export \\
\texttt{templates/conferences/} & HTML templates for all views \\
\texttt{docker-compose.yml} & Container orchestration for local/prod use \\
\bottomrule
\end{tabular}
\end{center}

% -------------------------------------------------------------------------
\section{Use of AI-Assisted Development}

The Spring 2026 team developed the entire codebase using
\textbf{Claude Code} (Anthropic, \texttt{claude-sonnet-4-6}) as a
development tool. This includes some views and exports, the organization
URL crawler, and this document.

Future teams should feel free to continue this approach. Claude Code can
read the existing codebase and extend it, provided the developer can articulate
the goal clearly. The \texttt{.tex} writeups in \texttt{ProjectFiles/}
document methodology and prompts used, and can serve as a model for
documenting AI-assisted work in future research outputs.

\end{document}
